% Source: https://artofproblemsolving.com/community/c3h3129215_2023_midmichigan_mathematical_olympiad__grades_79
% Stable identity: aops-topic-3129215-q04
(a) Use the vertices of a regular pentagon. Among its vertices, every distance is either a side length or a diagonal length. Any triangle therefore has three sides of only two possible lengths, so two sides are equal. No three pentagon vertices are collinear. (b) Add the centre of the pentagon. A triple containing the centre has two equal radii, while triples of vertices are covered by part (a). No line through two pentagon vertices contains the centre, because a regular pentagon has no antipodal pair. Hence no three of the six points are collinear.
\begin{asy}
usepackage("amsmath");
size(7cm);
pair[] V;for(int i=0;i<5;++i){V.push(dir(90+72*i));dot(V[i]);}draw(V[0]--V[1]--V[2]--V[3]--V[4]--cycle);dot((0,0));label("$O$",(0,0),S);draw((0,0)--V[0],dashed);draw((0,0)--V[2],dashed);
\end{asy}
